\clearpage
\fancyhead[L]{\small JUNO.}
\hypertarget{native-255-lower}{}
\begin{center}
\textbf{ON A NEW AND MOST REMARKABLE WANDERING STAR DISCOVERED BY INSPECTOR HARDING IN LILIENTHAL.}
\end{center}

\begin{center}
Monatliche Correspondenz. Vol. X. pp. 371–385 October 1804.
\end{center}

\begin{center}
— — —
\end{center}

I had scarcely sent the tireless and incomparable Dr. \textsc{Gauss} my first three observations of this celestial body, made on 13, 14, and 15 September, when I received the following reply by return post on 23 September. It will astonish all our readers as much as it astonished me. “What will you say,” writes this profound geometer, “when you hear that I dared to take my own observations, together with the three you were kind enough to send me and a couple of earlier ones by Dr. \textsc{Olbers}, which together cover only 14 days and a heliocentric arc of four degrees, that I dared, on so slippery a footing, to base a preliminary attempt and elliptic elements of an orbit on them without any hypothetical assumptions? The result can only be very \textit{precarious}. Still, I am inclined to hope that it can no longer depart \textit{enormously} or \textit{entirely} from the truth, but already gives at least a rough idea of the dimensions of the orbit. I flatter myself with still greater confidence that it will be enough to guide the search for the planet for perhaps a month, or even longer. And I can say with certainty that it represents all the observations made so far well. Here is the result for now. A comparison with the observations will follow shortly, at which point I shall also decide whether I can declare it worthy of the title \textit{Elements I of \textsc{Harding}’s Planet}.

\begin{center}
\begingroup\fontsize{9}{11}\selectfont\setlength{\tabcolsep}{3pt}
\begin{tabular}{lr}
Epoch, Seeberg meridian, 1804, 5 Sept.  & $24^\circ\ 53'\ 44''$ \\
Aphelion . . . . . . . . . . & $244\ 51\ 36$ \\
Ascending node . . . . . . & $171\ 48\ 24$ \\
Eccentricity . . . . . . . . . & $0.313757$ \\
Semimajor axis . . . . . . . . & $2.88208$ \\
Daily motion . . . . . . . & $725''18$ \\
Inclination of the orbit . . . . . . . & $15^\circ\ 12'\ 39''$ \\
Motion . . . . . . . . . . & direct \\
\end{tabular}\endgroup
\end{center}

“What do you say to this unusual orbit, its great eccentricity, and its close agreement with Ceres and Pallas as regards the axis and mean motion, \textit{an agreement that could very easily become exact with a slight change in the observation}? But I intend to remain true to my principle of mixing absolutely nothing hypothetical into the calculations and of not anticipating what future experi%
\SourceLeaf{256}%
ence may show. In a very short time we shall be much further along. There can now scarcely be any doubt that the orbit is worlds apart from a parabola and that \textsc{Harding}’s star deserves to be called a planet. It would therefore be desirable to give it a name soon. Of course, the right to christen it must be reserved for the discoverer alone, etc.

\Needspace{12\baselineskip}
Here are the observations Dr. \textsc{Gauss} made in Brunswick.

\begin{center}
\begingroup\fontsize{9}{11}\selectfont\setlength{\tabcolsep}{3pt}\renewcommand{\arraystretch}{1.12}
\begin{tabular}{r|r|r|r}
1804 & \shortstack{Mean time in\\Brunswick} & \shortstack{Apparent right\\ascension} & \shortstack{Apparent southern\\declination} \\ \hline
12 Sept. & $10^{u}\,35^{m}\,2^{s}$ & $0^\circ\ 54'\ 26''$ & $1^\circ\ 38'\ 15''$ \\
13 & $9\ 41\ 32$ & $0\ 45\ 24$ & $1\ 50\ 59$ \\
14 & $11\ 38\ 49$ & $0\ 35\ 37$ & $2\ 5\ 1$ \\
15 & $10\ 16\ 17$ & $0\ 26\ 53$ & $2\ 17\ 42$ \\
16 & $10\ 37\ 4$ & $0\ 17\ 17$ & $2\ 31\ 20$ \\
17 & $11\ 28\ 59$ & $0\ 7\ 23$ & $2\ 44\ 7$ \\
\end{tabular}\endgroup
\end{center}

\begin{center}
\begingroup\fontsize{9}{11}\selectfont\setlength{\tabcolsep}{3pt}\renewcommand{\arraystretch}{1.12}
\begin{tabular}{r|r|r|r}
1804 & \shortstack{Mean time in\\Brunswick} & \shortstack{Apparent right\\ascension} & \shortstack{Apparent southern\\declination} \\ \hline
18 & $11^{u}\,22^{m}\,16^{s}$ & $359^\circ\ 57'\ 26''$ & $2^\circ\ 58'\ 3''$ \\
21 & $10\ 24\ 52$ & $359\ 28\ 6$ & $3\ 37\ 59$ \\
24 & $10\ 1\ 2$ & $358\ 57\ 53$ & $4\ 18\ 6$ \\
25 & $8\ 44\ 25$ & $358\ 48\ 12$ & $4\ 30\ 44$ \\
27 & $10\ 20\ 29$ & $358\ 27\ 20$ & $4\ 57\ 47$ \\
28 & $8\ 29\ 4$ & $358\ 18\ 20$ & $5\ 10\ 22$ \\
\end{tabular}\endgroup
\end{center}

He writes about them: “The first three observations were made with a poor achromatic telescope, which was particularly poorly mounted, and therefore deserve little confidence. The subsequent ones, however, were made with a very good reflecting telescope. They will therefore be better, at least as good as the ring-micrometer method and the acuity of my eyesight allow.”

Just a few days after this letter, on 30 September, I received further confirmation of the elements of this planet’s orbit. “I am sending you,” writes Dr. \textsc{Gauss}, “new and improved elements. I flatter myself that they can already give an approximate determination of the true elements, and I therefore venture to announce them as the \textit{Elements I}.

\begin{center}
\begingroup\fontsize{9}{11}\selectfont\setlength{\tabcolsep}{3pt}
\begin{tabular}{lr}
Epoch, Seeberg meridian, 1804, 5 Sept. . . . . . . . & $20^\circ\ 38'\ 56''$ \\
Aphelion . . . . . . . . . . . . . . . . & $239\ 14\ 2$ \\
Ascending node . . . . . . . . . . . . . & $171\ 15\ 35$ \\
Eccentricity . . . . . . . . . . . . . . . . & $0.287359$ \\
Semimajor axis (logarithm) . . . . . . . . . . & $0.438682$ \\
Daily motion . . . . . . . . . . . . . . & $779''80$ \\
Inclination of the orbit . . . . . . . . . . . . . . & $13^\circ\ 34'\ 59''$ \\
Motion . . . . . . . . . . . . . . . . . & direct. \\
\end{tabular}\endgroup
\end{center}

Dr. \textsc{Gauss} then compared all the observations from Seeberg, Bremen, and Brunswick with these new elements. The differences are no greater than the errors possible with ring micrometers, and may in part also be due to the comparison stars.

\Needspace{13\baselineskip}
\begin{center}
\textit{Seeberg observations:}
\end{center}

\begin{center}
\begingroup\fontsize{8}{10}\selectfont\setlength{\tabcolsep}{3pt}\renewcommand{\arraystretch}{1.12}
\begin{tabular}{r|r|r|r|r|r}
1804 & \shortstack{Mean time\\at Seeberg} & \shortstack{Calculated right\\ascension} & \shortstack{Difference\\R—B} & \shortstack{Calculated\\declination} & \shortstack{Diff.\\R—B} \\ \hline
13 Sept. & $12^{u}\,31^{m}\,59^{s}$ & $0^\circ\ 44'\ 57''5$ & $+1''0$ & $1^\circ\ 52'\ 34''7$ & $-2''3$ \\
14 & $12\ 27\ 27$ & $0\ 35\ 43.8$ & $+2.8$ & $2\ 5\ 35.3$ & $-0.3$ \\
15 & $12\ 22\ 53$ & $0\ 26\ 20.4$ & $+0.4$ & $2\ 18\ 40.5$ & $+2.0$ \\
17 & $12\ 13\ 45$ & $0\ 7\ 7.8$ & $-0.3$ & $2\ 45\ 8.7$ & $-4.4$ \\
18 & $12\ 9\ 10$ & $359\ 57\ 20.4$ & $0$ & $2\ 58\ 28.3$ & $-2.6$ \\
20 & $ .\quad .\quad . $ & $359\ 37\ 28.7$ & $ .\ .\ . $ & $3\ 25\ 14.9$ & $+11.4$ \\
23 & $11\ 46\ 11$ & $359\ 7\ 11.9$ & $-3.0$ & $4\ 5\ 32.0$ & $-2.5$ \\
\end{tabular}\endgroup
\end{center}
\Needspace{23\baselineskip}
\SourceLeaf{257}%
\begin{center}
\textit{Bremen observations:}
\end{center}

\begin{center}
\begingroup\fontsize{8}{10}\selectfont\setlength{\tabcolsep}{3pt}\renewcommand{\arraystretch}{1.12}
\begin{tabular}{r|r|r|r|r|r}
1804 & \shortstack{Mean time\\in Bremen} & \shortstack{Calculated right\\ascension} & \shortstack{Diff.\\R—B} & \shortstack{Calculated\\declination} & \shortstack{Diff.\\R—B} \\ \hline
7 Sept. & $10^{u}\,37^{m}\,21^{s}$ & $1^\circ\ 36'\ 57''$ & $+5''$ & $0^\circ\ 36'\ 6''$ & $-3''$ \\
8 & $8\ 11\ 20$ & $1\ 29\ 42$ & $+16$ & $0\ 47\ 4$ & $-15$ \\
9 & $10\ 48\ 50$ & $1\ 20\ 36$ & $+10$ & $1\ 0\ 47$ & $-3$ \\
10 & $8\ 15\ 6$ & $1\ 13\ 5$ & $+5$ & $1\ 11\ 58$ & $+2$ \\
11 & $10\ 34\ 3$ & $1\ 3\ 35$ & $+11$ & $1\ 25\ 51$ & $+10$ \\
12 & $11\ 18\ 32$ & $0\ 54\ 27$ & $+22$ & $1\ 39\ 3$ & $-1$ \\
13 & $8\ 54\ 0$ & $0\ 46\ 17$ & $+14$ & $1\ 50\ 41$ & $-9$ \\
14 & $8\ 24\ 44$ & $0\ 37\ 15$ & $+8$ & $2\ 3\ 27$ & $ .\quad . $ \\
15 & $10\ 54\ 28$ & $0\ 26\ 52$ & $+12$ & $2\ 17\ 57$ & $-8$ \\
17 & $10\ 23\ 9$ & $0\ 7\ 50$ & $+25$ & $2\ 44\ 12$ & $-20$ \\
18 & $8\ 38\ 17$ & $359\ 58\ 44$ & $-3$ & $2\ 56\ 35$ & $-16$ \\
21 & $8\ 30\ 54$ & $359\ 28\ 49$ & $-3$ & $3\ 36\ 50$ & $-4$ \\
21 & $10\ 9\ 32$ & $359\ 28\ 8$ & $-1$ & $3\ 37\ 45$ & $-1$ \\
23 & $13\ 25\ 57$ & $359\ 6\ 26$ & $+8$ & $4\ 6\ 32$ & $-5$ \\
24 & $8\ 27\ 37$ & $358\ 58\ 22$ & $+8$ & $4\ 17\ 12$ & $-10$ \\
25 & $8\ 41\ 38$ & $358\ 48\ 15$ & $+2$ & $4\ 30\ 32$ & $-22$ \\
\end{tabular}\endgroup
\end{center}

\begin{center}
\textit{Brunswick observations;} 
\end{center}

\begin{center}
\begingroup\fontsize{8}{10}\selectfont\setlength{\tabcolsep}{3pt}\renewcommand{\arraystretch}{1.12}
\begin{tabular}{r|r|r|r|r|r}
1804 & \shortstack{Mean time in\\Brunswick} & \shortstack{Calculated right\\ascension} & \shortstack{Diff.\\R—B} & \shortstack{Calculated\\declination} & \shortstack{Diff.\\R—B} \\ \hline
12 Sept. & $10^{u}\,35^{m}\,2^{s}$ & $0^\circ\ 54'\ 46''$ & $+15''$ & $1^\circ\ 38'\ 36''$ & $+21''$ \\
13 & $9\ 41\ 32$ & $0\ 46\ 2$ & $+38$ & $1\ 51\ 2$ & $+3$ \\
14 & $11\ 38\ 49$ & $0\ 35\ 58$ & $+21$ & $2\ 5\ 15$ & $+14$ \\
15 & $10\ 16\ 17$ & $0\ 27\ 10$ & $+17$ & $2\ 17\ 33$ & $-2$ \\
16 & $10\ 37\ 4$ & $0\ 17\ 28$ & $+11$ & $2\ 30\ 58$ & $-22$ \\
17 & $11\ 28\ 59$ & $0\ 7\ 25$ & $+2$ & $2\ 44\ 44$ & $+15$ \\
18 & $11\ 22\ 16$ & $359\ 57\ 38$ & $+12$ & $2\ 58\ 4$ & $+1$ \\
21 & $10\ 24\ 52$ & $359\ 28\ 5$ & $-1$ & $3\ 37\ 50$ & $-10$ \\
24 & $10\ 1\ 2$ & $358\ 57\ 45$ & $-8$ & $4\ 18\ 1$ & $-5$ \\
25 & $8\ 44\ 25$ & $358\ 48\ 6$ & $-6$ & $4\ 30\ 43$ & $-1$ \\
27 & $10\ 29\ 29$ & $358\ 27\ 8$ & $-12$ & $4\ 58\ 17$ & $+30$ \\
28 & $8\ 29\ 4$ & $358\ 17\ 50$ & $-30$ & $5\ 10\ 30$ & $+8$ \\
\end{tabular}\endgroup
\end{center}

Dr. \textsc{Olbers} was also kind enough to send us his further observations of the remarkable stranger:

\begin{center}
\begingroup\fontsize{9}{11}\selectfont\setlength{\tabcolsep}{3pt}\renewcommand{\arraystretch}{1.12}
\begin{tabular}{r|r|r|r}
1804 & \shortstack{Mean time\\in\\Bremen} & \shortstack{Apparent\\right\\ascension} & \shortstack{Apparent\\southern\\declination} \\ \hline
9 Sept. & $10^{u}\,48^{m}\,50^{s}$ & $1^\circ\ 20'\ 30''$ & $1^\circ\ 1'\ 5''$ \\
10 & $8\ 15\ 6$ & $1\ 12\ 55$ & $1\ 11\ 55$ \\
11 & $10\ 43\ 54$ & $1\ 3\ 20$ & $1\ 25\ 48$ \\
\end{tabular}\endgroup
\end{center}

Dr. \textsc{Olbers} writes about them: “These observations still need a small correction, since all the comparisons, with a single exception, had to be made with faint stars from the \textit{Hist. cél.} As \textsc{Harding}’s observation and my own are probably the first made of this new planet, it would be very desirable if you would be so kind as to determine these faint stars of \textsc{La Lande} more accurately. The stars are as follows:

\Needspace{11\baselineskip}
\begin{center}\fontsize{9}{11}\selectfont
\begin{tabular}{l}
p. 119 Hist. cél.\\
Magnitude $8$\quad third wire $0^u\ 8^m\ 53^s5$\\
\quad (should be $0^u\ 7^m\ 53^s5$)\quad zenith distance $49^\circ\ 20'\ 6''$
\end{tabular}
\[
\text{p. 131 Hist. cél.}\left\{\begin{array}{lll}
8\ \text{—}&\text{middle wire}\ 0\ 3\ 40.5&\text{— —}\ 49\ 59\ 0\\
8\ \text{—}&\text{— —}\ 0\ 4\ 26.0&\text{— —}\ 50\ 15\ 12
\end{array}\right.
\]
\end{center}
\par
\SourceLeaf{258}%
\indent I too determined these three stars as accurately as possible. Here are their mean positions for the beginning of 1804:

\begin{center}
\begingroup\fontsize{8}{10}\selectfont\setlength{\tabcolsep}{3pt}\renewcommand{\arraystretch}{1.12}
\begin{tabular}{r|r|r|r|r}
Magnitude & \shortstack{Mean right\\ascension\\1804} & \shortstack{Annual\\change} & \shortstack{Mean southern\\declination\\1804} & \shortstack{Annual\\change} \\ \hline
$8$ & $1^\circ\ 0'\ 24''0$ & $+46''0$ & $1^\circ\ 7'\ 19''6$ & $-20''0$ \\
$8$ & $1\ 11\ 52.7$ & $+46.0$ & $1\ 23\ 30.9$ & $-20.0$ \\
$8$ & $2\ 0\ 20.9$ & $+46.0$ & $0\ 28\ 15.7$ & $-20.0$ \\
\end{tabular}\endgroup
\end{center}

It is unprecedented in the history of astronomy in any age or nation, and bears witness to the brilliant era of astronomy in Germany today, that a planet should have been predicted and, in the short space of three weeks, discovered and observed, its orbit calculated, and its future course plotted. Yet all this was done by the combined efforts of four German astronomers, who had accomplished it all before news of the existence of this new celestial body had even reached our jealous neighbours.

Finally, here is an ephemeris of the future geocentric course of this new planet, which Dr. \textsc{Gauss} calculated from his Elements I above. However much improvement these elements may still need, Dr. \textsc{Gauss} confidently hopes that they will accurately represent this course for several weeks. The following ephemeris will therefore be very welcome to distant astronomers who have yet to find this new guest in the heavens. The times are calculated for midnight at Seeberg.

\begin{center}
\begingroup\fontsize{9}{11}\selectfont\setlength{\tabcolsep}{3pt}\renewcommand{\arraystretch}{1.12}
\begin{tabular}{r|r|r}
\shortstack{Midnight\\at Seeberg} & \shortstack{Right ascension\\of the new\\planet} & \shortstack{Southern declination\\of the new\\planet.} \\ \hline
5 Sept. & $1^\circ\ 52'$ & $0^\circ\ 12'$ \\
12 & $0\ 54$ & $1\ 39$ \\
19 & $359\ 47$ & $3\ 12$ \\
26 & $358\ 36$ & $4\ 55$ \\
30 & $357\ 56$ & $5\ 39$ \\
3 Oct. & $357\ 27$ & $6\ 17$ \\
6 & $357\ 0$ & $6\ 54$ \\
9 & $356\ 35$ & $7\ 29$ \\
12 & $356\ 12$ & $8\ 2$ \\
15 & $355\ 51$ & $8\ 33$ \\
\end{tabular}\endgroup
\end{center}
\begin{center}
\begingroup\fontsize{9}{11}\selectfont\setlength{\tabcolsep}{3pt}\renewcommand{\arraystretch}{1.12}
\begin{tabular}{r|r|r}
\shortstack{Midnight\\at Seeberg} & \shortstack{Right ascension\\of the new\\planet} & \shortstack{Southern declination\\of the new\\planet} \\ \hline
18 Oct. & $355^\circ\ 35'$ & $9^\circ\ 1'$ \\
21 & $355\ 21$ & $9\ 27$ \\
24 & $355\ 11$ & $9\ 50$ \\
27 & $355\ 4$ & $10\ 11$ \\
30 & $355\ 2$ & $10\ 28$ \\
2 Nov. & $355\ 4$ & $10\ 42$ \\
5 & $355\ 10$ & $10\ 54$ \\
8 & $355\ 20$ & $11\ 3$ \\
11 & $355\ 34$ & $11\ 8$ \\
14 & $355\ 53$ & $11\ 13$ \\
\end{tabular}\endgroup
\end{center}

\hfill[\textsc{von Zach.}]
